\section{第二章 城市轨道交通变电所高压开关柜在线监测技术研究}

\section{高压开关柜故障分析与监测需求}

（一）

城市轨道交通变电所高压开关柜作为变电系统的核心电气设备，其结构设计和工作原理直接决定了设备的运行特性和故障特征。目前城市轨道交通变电所主要采用两种形式的开关柜：金属封闭开关设备（GIS）和敞开式开关设备。其中GIS设备由于其占地面积小、可靠性高、维护量少等显著优点，在城市轨道交通领域得到了广泛应用。GIS设备采用全封闭结构，将所有带电部件封闭在金属外壳内，并用SF6气体作为绝缘和灭弧介质，这种设计不仅提高了设备的安全性，还显著减少了设备的占地空间，特别适合城市轨道交通变电所这种空间受限的应用环境。

高压开关柜的结构较为复杂，主要由断路器、隔离开关、互感器、避雷器和操动机构等核心部件组成。断路器作为开关柜的核心部件，主要负责实现电路的通断控制，它能够在正常工作状态和故障状态下可靠地接通和断开电路；隔离开关则主要提供电气隔离功能，在检修时确保检修人员的安全；互感器负责实现电压和电流的测量，为保护装置和测量仪表提供准确的电气参数；避雷器则提供过电压保护功能，防止雷击或其他过电压对设备造成损害；操动机构作为开关柜的执行部件，为断路器和隔离开关提供分合闸所需的动力和控制信号。这些部件相互配合，共同构成了高压开关柜完整的电气系统。

\begin{figure}[h]
\centering
%% % \includegraphics[width=0.9\textwidth]{figs/switchgear_structure.png}
\caption{高压开关柜结构组成示意图}
\label{fig:switchgear_structure}
\end{figure}

（一）

通过对城市轨道交通变电所高压开关柜长期运行数据的分析和故障案例研究，可以发现高压开关柜的故障类型多种多样，但主要可以归纳为绝缘故障、机械故障、温升故障和其他故障四大类。这些故障不仅具有不同的发生机理，而且对设备安全运行的影响程度也各不相同。

绝缘故障是高压开关柜最常见也是最严重的故障类型之一，主要包括固体绝缘故障、气体绝缘故障和液体绝缘故障三种形式。固体绝缘故障通常表现为绝缘子表面污秽和绝缘子老化等问题，这些问题往往是由长期运行积累的灰尘、潮气以及材料自然老化引起的；气体绝缘故障主要发生在SF6气体绝缘的GIS设备中，表现为SF6气体泄漏和气体绝缘性能下降，这些问题可能由密封圈老化、焊接质量缺陷或外部碰撞造成；液体绝缘故障则主要涉及绝缘油的劣化和油中水分超标等问题，这些故障通常与油的氧化变质和水分侵入有关。

机械故障主要涉及开关柜的操动机构和传动系统，是影响设备可靠性的重要因素。操动机构卡涩往往是由于润滑油不足、机械零件磨损或异物进入造成的，这种故障会导致分合闸操作困难，严重时可能造成拒动；分合闸失败则是机械故障的典型表现形式，通常由储能弹簧失效、电磁铁故障或传动部件损坏引起；触头烧损问题则主要发生在分合闸过程中，由于电弧效应导致触头表面烧伤，这种故障不仅影响导电性能，还可能引发连锁故障。

温升故障作为一种常见的运行状态异常，主要由多种因素共同作用引起。接触电阻增大是温升故障的主要原因之一，当触头表面氧化或出现不平整现象时，会导致接触电阻显著增加，进而产生过多的热量；负荷电流过大也是造成温升故障的重要因素，当实际运行电流超过设备额定电流时，会产生过量的热能；环境温度过高同样会加剧温升问题，特别是设备安装在通风不良的封闭空间时，更容易积累热量。

除此之外，高压开关柜还可能发生其他类型的故障，如辅助回路故障、密封故障和接地故障等。这些故障虽然发生频率相对较低，但同样需要引起足够的重视。辅助回路故障主要表现为控制回路断开和信号回路异常等问题；密封故障主要涉及SF6气体密封圈的老化和压力异常；接地故障则表现为接地线松动和接地电阻增大等问题。

\begin{figure}[h]
\centering
%% % \includegraphics[width=0.8\textwidth]{figs/fault_mechanism.png}
\caption{高压开关柜故障机理分析图}
\label{fig:fault_mechanism}
\end{figure}

（一）

基于对高压开关柜故障类型和机理的深入分析，可以明确在线监测系统需要满足的基本需求和技术要求。这些需求不仅涵盖了监测参数的全面性，还包括监测精度、实时性和可靠性等多个维度，是构建有效在线监测系统的基本保障。

在监测参数的全面性方面，高压开关柜在线监测系统需要覆盖电气参数、机械参数、热学参数和绝缘参数等多个方面。电气参数主要包括电压、电流、功率等基础电气量，这些参数反映了设备的基本运行状态；机械参数主要涉及行程、速度、力等机械特性指标，这些参数能够反映操动机构的运行状况；热学参数则包括温度、温升率等热学特征，这些参数对于发现温升故障具有重要意义；绝缘参数主要包括局放、绝缘电阻等绝缘特性指标，这些参数能够提前发现绝缘故障隐患。

监测精度的要求直接决定了在线监测系统的诊断准确性。对于温度监测，系统需要达到±1℃的监测精度，以确保能够敏感地捕捉温度异常变化；对于局放监测，系统需要具备1pC的监测灵敏度，才能有效检测微弱的局部放电信号；对于机械特性监测，系统需要达到±0.1mm的监测精度，以准确反映触头的位移变化。这些精度要求是确保故障诊断准确性的基本保障。

实时性要求是评价在线监测系统性能的重要指标。数据采集频率不应低于100Hz，以确保能够捕捉到快速变化的电气和机械参数；数据传输延迟应小于1秒，以保证监测数据的时效性；报警响应时间应小于10秒，以确保能够在故障发生初期及时发出预警信号。这些实时性要求对于实现故障的早期发现和及时处理具有重要意义。

可靠性要求则是确保在线监测系统长期稳定运行的关键。系统可用率应大于99.9\%，以保证监测系统的连续运行能力；误报率应小于1\%，避免频繁的虚假报警影响维护人员的判断；漏报率应小于0.1\%，确保重要的故障信息不会被遗漏。这些可靠性要求对于提高整个监测系统的实用价值具有决定性意义。

\begin{figure}[h]
\centering
%% % \includegraphics[width=0.8\textwidth]{figs/monitoring_requirements.png}
\caption{在线监测系统技术需求分析}
\label{fig:monitoring_requirements}
\end{figure}

\section{在线监测技术实现}

（一）

温度监测作为高压开关柜在线监测的基础技术，对于发现温升故障和预防热故障具有重要意义。目前主流的温度监测技术主要包括红外测温技术和光纤测温技术两种类型，每种技术都具有独特的优势和应用特点。

红外测温技术利用红外辐射原理实现非接触式温度测量，具有诸多显著优势。该技术最大的特点是非接触测量，不会影响设备的正常运行，能够在设备带电状态下进行监测；响应速度快，能够快速捕捉温度变化，特别适用于动态监测场景；测量范围广，可以同时监测多个测点，提高监测效率。红外测温系统的组成相对复杂，主要包括红外热像仪、数据采集单元和图像处理软件等部分。红外热像仪负责采集温度场分布信息，能够生成直观的温度分布图像；数据采集单元负责处理红外信号，将模拟信号转换为数字信号；图像处理软件则对温度数据进行深入分析，包括温度趋势分析、热点识别等功能。

光纤测温技术作为一种新兴的温度监测手段，基于光纤布拉格光栅（FBG）原理，具有传统测温技术无法比拟的优势。该技术最大的特点是抗电磁干扰能力强，能够在强电磁场环境下稳定工作，这对于高压开关柜这种强电设备特别重要；支持多点分布式测量，能够同时监测多个位置的温度变化；具有高精度和高可靠性的特点，能够提供准确的温度测量结果。光纤测温系统的实现方案主要包括光纤传感器的布设、光谱分析仪和数据处理单元三个核心部分。光纤传感器通常沿母线和触头等关键部位布置，形成分布式监测网络；光谱分析仪负责解调温度信息，将光信号转换为温度数据；数据处理单元则负责温度补偿和校正，确保测量结果的准确性。

\begin{figure}[h]
\centering
%% % \includegraphics[width=0.8\textwidth]{figs/temperature_monitoring.png}
\caption{温度监测技术对比分析}
\label{fig:temperature_monitoring}
\end{figure}

（一）

局部放电监测是发现绝缘故障的重要手段，能够提前发现绝缘劣化问题。目前常用的局放监测技术主要包括超声波法、高频法和特高频法三种，每种方法都有其独特的检测原理和应用优势。

超声波法局放监测利用局部放电产生的超声波信号进行检测，具有结构简单、成本较低的特点。该方法的实现主要包括传感器布置、信号处理和定位算法三个方面。传感器通常采用压电传感器，安装在开关柜外壳表面，能够接收局部放电产生的超声波信号；信号处理环节包括信号放大、滤波和特征提取等步骤，通过数字信号处理技术提取有用的放电特征；定位算法则基于到达时间差原理，通过多个传感器的时差测量确定放电源的具体位置。

高频法局放监测通过检测局部放电产生的高频电流脉冲来识别放电现象，具有检测灵敏度高的特点。该方法主要包括耦合电容器、高频CT和信号分析三个组成部分。耦合电容器负责提取高频信号，能够将高频放电信号从主回路中分离出来；高频CT则专门检测接地线上的高频电流变化；信号分析环节主要进行频谱分析和模式识别，通过分析信号的频谱特征和时域特征来识别不同的放电类型。

特高频法局放监测通过检测局部放电产生的电磁波辐射来发现绝缘故障，具有定位精度高的优势。该方法主要由天线传感器、放大器和频谱仪三个部分组成。天线传感器负责接收特高频电磁波信号，能够捕获局部放电辐射的电磁波；放大器负责信号放大和滤波，提高检测灵敏度；频谱仪则对信号进行频谱分析和特征提取，通过分析电磁波的频谱特征来判断放电的严重程度。

\begin{figure}[h]
\centering
%% % \includegraphics[width=0.9\textwidth]{figs/pd_monitoring_methods.png}
\caption{局放监测技术方法对比}
\label{fig:pd_monitoring_methods}
\end{figure}

（一）

机械特性监测主要针对开关柜的分合闸性能进行监测，是发现机械故障的重要手段。目前主要的机械特性监测技术包括行程特性监测、速度特性监测和力特性监测三种类型，共同构成了完整的机械状态评估体系。

行程特性监测通过位移传感器测量动触头的行程变化，能够直接反映分合闸过程的机械特性。该监测方法主要包括线性位移传感器、数据采集和特性分析三个环节。线性位移传感器安装在操动机构上，能够精确测量触头的位移变化；数据采集系统需要同步采集电压、电流等电气信号，以便进行综合分析；特性分析则主要进行行程-时间曲线的分析，通过分析曲线的形状和特征来判断机械状态是否正常。

速度特性监测通过对行程数据的微分计算获得触头的运动速度，能够发现机械运动的异常变化。该方法需要高速数据采集系统，采样频率不低于10kHz，以确保能够捕捉到快速的机械运动；数字滤波环节负责去除高频噪声，提高数据质量；速度计算则采用数值微分算法，将位移数据转换为速度数据。

力特性监测主要测量操动力和触头接触压力，能够反映机械传动的受力情况。该监测方法主要采用力传感器和压力传感器两种类型的传感器。力传感器安装在操动机构的关键部位，能够测量分合闸过程中的力变化；压力传感器则专门测量触头接触压力，确保接触的可靠性；动态分析环节则进行力-时间曲线的分析，通过分析力曲线的特征来判断机械状态。

\begin{figure}[h]
\centering
%% % \includegraphics[width=0.8\textwidth]{figs/mechanical_monitoring.png}
\caption{机械特性监测系统架构}
\label{fig:mechanical_monitoring}
\end{figure}

\section{多参数综合监测系统设计}

（一）

多参数综合监测系统的设计需要充分考虑高压开关柜的运行特点和技术需求，采用模块化、分层式的架构设计理念，确保系统具有良好的可扩展性、可维护性和可靠性。整个系统架构从硬件和软件两个维度进行设计，形成完整的监测体系。

硬件架构作为系统的基础平台，主要包括传感器层、数据采集层、通信层和控制层四个层次。传感器层是系统的数据来源，主要部署各类监测传感器，包括温度传感器、局放传感器、机械特性传感器等，这些传感器负责采集设备的各种运行参数；数据采集层负责信号调理和A/D转换，将传感器采集的模拟信号转换为数字信号，并进行初步的信号处理；通信层负责数据传输和网络通信，采用工业以太网等可靠的通信方式确保数据传输的实时性和可靠性；控制层则由主控单元和嵌入式处理器组成，负责整个系统的协调控制和数据管理。

软件架构采用分层设计理念，将复杂的系统功能分解为多个相对独立的层次，便于系统的开发和维护。数据采集层负责实时数据采集和预处理，包括数据采集调度、数据校验和数据预处理等功能；数据处理层主要进行信号分析和特征提取，通过数字信号处理技术提取有用的特征信息；诊断决策层负责故障诊断和预测，基于特征信息进行智能分析和决策；人机界面层则提供监控显示和操作管理功能，为用户提供友好的操作界面和丰富的功能选项。

\begin{figure}[h]
\centering
%% % \includegraphics[width=0.9\textwidth]{figs/system_architecture.png}
\caption{多参数综合监测系统架构设计}
\label{fig:system_architecture}
\end{figure}

（一）

多源数据融合技术是提高监测系统诊断准确性的关键技术，通过综合利用多种传感器的数据，可以获得更加全面和准确的设备状态信息。数据融合技术按照融合层次的不同可以分为数据级融合、特征级融合和决策级融合三种类型。

数据级融合是最基础的融合方式，主要对原始数据进行直接融合处理。这种方法能够充分利用所有传感器的原始信息，但对数据处理能力要求较高；特征级融合则对各传感器提取的特征向量进行融合，能够减少数据量，提高处理效率；决策级融合则对各传感器的诊断结果进行融合，通过投票或加权方式得到最终的诊断结论，具有较强的容错能力。

基于大数据分析技术的故障分析方法为设备状态评估提供了新的思路和手段。统计分析方法通过对历史数据的统计特征分析，发现设备运行的规律和异常模式；模式识别技术则通过机器学习算法对故障模式进行分类，能够自动识别不同类型的故障；预测建模则基于机器学习算法建立故障预测模型，能够提前预测故障发生的概率和时间。

\begin{figure}[h]
\centering
%% % \includegraphics[width=0.8\textwidth]{figs/data_fusion_process.png}
\caption{多源数据融合处理流程}
\label{fig:data_fusion_process}
\end{figure}

（一）

故障预警模型是监测系统的重要组成部分，能够根据设备状态信息预测潜在故障风险，为维护决策提供重要依据。目前常用的预警模型主要包括阈值预警模型和智能预警模型两种类型。

阈值预警模型是一种相对简单的预警方法，通过设置固定阈值来判断设备状态是否异常。该模型包括单参数阈值、多参数综合和趋势分析三种形式。单参数阈值主要针对温度、局放等单一指标设置阈值，当超过阈值时发出报警；多参数综合则采用加权组合的方式，对多个参数进行综合评价；趋势分析则基于参数的变化率进行预警，能够发现渐变式的故障发展。

智能预警模型利用人工智能技术实现更加准确和智能的故障预警，具有较强的学习和适应能力。该模型主要采用支持向量机（SVM）、神经网络和深度学习等人工智能算法。支持向量机特别适用于故障分类问题，能够有效区分正常状态和故障状态；神经网络则擅长处理非线性预测问题，能够建立复杂的故障预测模型；深度学习技术则通过多层神经网络实现复杂模式识别，能够发现传统方法难以识别的故障模式。

\begin{figure}[h]
\centering
%% % \includegraphics[width=0.8\textwidth]{figs/warning_model_comparison.png}
\caption{故障预警模型对比分析}
\label{fig:warning_model_comparison}
\end{figure}